\chapter{Gauge Field Fluid Dynamics}
\label{chap:gauge-fluid}

\section{Fluid Dynamics of the Brane}
\label{sec:brane-fluid}

\begin{itemize}
  \item Describe the fluid dynamics on the brane whose thickness is determined by
\end{itemize}
  the substrate parameter $q$ (\hyperref[def:substrate-parameter]{Definition}).
\begin{itemize}
  \item TODO: investigate whether the viscosity gradient can be encoded by the
\end{itemize}
  $q$-Laplacian $\Delta_q$.
\begin{itemize}
  \item TODO: formulate a Navier--Stokes analogue in the $q$-deformed gauge-field
\end{itemize}
  setting.

\section{Geodesics as Stable Vortex Points}
\label{sec:geodesic-vortices}

\begin{itemize}
  \item Geodesics on $M$ as stable fixed points of vortex configurations in the
\end{itemize}
  brane fluid.
\begin{itemize}
  \item TODO: perform stability analysis via the $q$-deformed energy functional.
\end{itemize}

\section{Braid Operators in the Fluid Picture}
\label{sec:braid-fluid}

TODO: work out the connection between braid-group operators
(\hyperref[chap:path-integrals-braiding]{Chapter}) and vortex braiding in the fluid
dynamics.

\section{$q$-Parameter, Planck Limit, and Reynolds Number}
\label{sec:planck-reynolds}

TODO: relate the $q$-deformation parameter and the Planck limit to the
critical separation distance and the Reynolds number of the brane fluid.
